% !TEX root = ../main.tex
\chapter{Follow the winner}

The first principle in the book is called follow the winner \parencite[Ch.~4]{lihoi}: the
algorithms increase the weight of the assets that have done well so far, they "follow the winner". 
The objective is to do almost as well as the the best constant portfolio in hindsight. 
The problem is that the BCRP needs future prices, so an online algorithm cannot be the
BCRP, it can only get close to it. To measure how close, we compare the final wealth of
both with the regret,
\[
\mathrm{regret}_n(\mathrm{Alg})=\log S_n(\mathrm{BCRP})-\log S_n(\mathrm{Alg})
=\log\frac{S_n(\mathrm{BCRP})}{S_n(\mathrm{Alg})}.
\]
This regret represents how many times richer the BCRP ended than the algorithm.
A regret of 0 means that both ended with the same wealth, and a regret of
0.1 means that the BCRP ended about 10\,\% richer. \\
An algorithm is called universal if $\mathrm{regret}_n/n\to 0$ \parencite{cover91}. 
This means that the difference in average daily growth between the algorithm and the
BCRP goes to zero as the number of days grows, whatever the prices do. 
This does not imply that the final wealth is equal, but that in the limit they grow at the same speed. 
The three algorithms of this chapter are universal.

\section{Universal portfolio}
The universal portfolio (UP) was proposed by \textcite{cover91}. Instead of choosing one
portfolio $b$, it follows all of them at the same time. Imagine one investor for each
possible portfolio $b\in\Delta_m$, all with the same money and each one following his own
CRP. After $t$ days, the wealth of each investor of $b$ is
\[
S_t(b)=\prod_{\tau=1}^{t}b\cdot x_\tau .
\]
Each day, UP takes the average of all the portfolios, but the richer investors count
more because their assets have increased in value:
\begin{equation}
b_{t+1}=\frac{\int_{\Delta_m}b\,S_t(b)\,\mathrm{d}\mu(b)}{\int_{\Delta_m}S_t(b)\,\mathrm{d}\mu(b)} .
\label{eq:up}
\end{equation}
The numerator adds up every portfolio $b$ multiplied by the wealth of its investor. The
denominator is the total wealth, and dividing by it makes the weights add up to one. We
use integrals instead of sums because there are infinitely many portfolios. The measure
$\mu$ says how the initial money is shared among the investors, and we use the uniform
one, which gives the same amount to all, as in the original paper. 
\\ The CRPs that have done well weigh more, so UP follows the winners. The final wealth
of UP is the average of the final wealth of all the CRPs. \textcite{cover91} proved that
the regret of UP is $O(m\log n)$, which grows very slowly with $n$, so UP is universal. \\

The problem is that these integrals are very expensive to compute when there are many
assets, because the work grows exponentially with $m$. Since that was the case the
integrals were simplified to sums. We choose $K=1000$ random portfolios $b^{(1)},\dots,b^{(K)}$ 
once, all with the same probability (a Dirichlet distribution with all parameters equal to
one), and use sums over them:
\[
b_{t+1}\approx\frac{\sum_{k=1}^{K}b^{(k)}\,S_t(b^{(k)})}{\sum_{k=1}^{K}S_t(b^{(k)})} .
\]
This simplification works badly when there are many assets and many days. UP is an average of our random 
portfolios, so it cannot be more concentrated than they are. If we split the money at random among a few 
assets, some of them get a large share, so there are portfolios with a strong bet on one asset. If we 
split it among many assets, the money is spread so thin that every asset gets a small share, and none 
of the portfolios bets heavily on a single one. Since UP is an average of them, it cannot concentrate on the
best asset either, and it stays close to the equal-weight CRP. % which is what happens in the DJIA universe


\section{Exponential gradient}

Exponential gradient (EG) was proposed by \textcite{helmbold98}. It does not use the
whole history, only the last day. After seeing day $t$, it looks for a new portfolio that would have done 
better than the current one that day, but that is not too different from the current one:
\begin{equation}
b_{t+1}=\arg\max_{b\in\Delta_m}\ \eta\,\log(b\cdot x_t)-R(b,b_t),
\qquad R(b,b_t)=\sum_{i=1}^m b_i\log\frac{b_i}{b_{t,i}} .
\label{eq:eg}
\end{equation}
The first term measures how well $b$ would have done on day $t$. The second term is called the
relative entropy, and acts as a regularization term measuring how different $b$ is from the current
portfolio $b_t$. The learning rate $\eta>0$ says how much we care about the last day compared
with staying where we are. The logarithm makes this problem hard to solve, so the
authors replace it by its first-order approximation around $b_t$, so that the solution can be written directly:
\begin{equation}
b_{t+1,i}=\frac{b_{t,i}\,\exp\!\big(\eta\,x_{t,i}/(b_t\cdot x_t)\big)}{Z},
\qquad i=1,\dots,m,
\label{eq:eg-update}
\end{equation}
where $Z$ normalizes the weights. The new weight is the old 
one multiplied by a factor. The ratio $x_{t,i}/(b_t\cdot x_t)$ compares the return of
asset $i$ with the return of the whole portfolio. If the asset did better than the
portfolio, the ratio is larger than one, the factor is larger than one and the weight
grows. If it did worse, the weight decreases. After many days, an asset that keeps
beating the portfolio keeps gaining weight, so EG follows the winners.

The learning rate controls the speed. If $\eta$ is very small the portfolio hardly moves
from the uniform one, so EG behaves like the equal-weight CRP, and if it is large the
portfolio reacts strongly to each single day. 

By default we use $\eta=0.05$, the value suggested by \textcite{helmbold98} and used in
the experiments of the book. In our experiments $\eta$ is either this default or chosen % reference walk-forward when mentioned
each year from a grid of values, using only past data. The theory asks for something
different. With a fixed $\eta$, the guarantee only says that the regret per day stays
below a constant, not that it goes to zero, so EG is not universal. If $\eta$ gets 
smaller as the number of days grows, of the order of % not implemented neither in the book nor in the code
$\sqrt{(\log m)/n}$, the regret of EG is $O(\sqrt{n\log m})$
\parencite{helmbold98}, and then EG is universal \parencite[Sec.~4.2]{lihoi}.


\section{Online Newton step}

Online Newton step (ONS) was proposed by \textcite{agarwal06}, who adapted to portfolio
selection the Newton method for online convex optimization \parencite{hazan07}. The
Newton method is used to find the minimum of a function. It uses the first derivative to
know in which direction to go and the second derivative to know how big the step should
be. ONS does the same with the logarithm of the wealth. For each day we define
\[
g_\tau=\frac{x_\tau}{b_\tau\cdot x_\tau},\qquad
A_t=I_m+\sum_{\tau=1}^{t}g_\tau g_\tau^\top,\qquad
c_t=\Big(1+\frac1\beta\Big)\sum_{\tau=1}^{t}g_\tau .
\]
The vector $g_\tau$ is the derivative of $\log(b\cdot x_\tau)$ with respect to $b$, at the
portfolio held that day. It points towards the portfolios that would have grown more. The 
matrix $g_\tau g_\tau^\top$ describes the curvature of $\log(b\cdot x_\tau)$, that is, its 
second derivative. So $c_t$ adds up all the derivatives seen so far, which gives the direction, 
and $A_t$ adds up all the second derivatives, which gives the curvature that sets the size of the
step. ONS starts from the uniform portfolio and then
\[
b_{t+1}=\Pi^{A_t}_{\Delta_m}\big(\delta\,A_t^{-1}c_t\big),
\qquad
\Pi^{A}_{\Delta_m}(y)=\arg\min_{b\in\Delta_m}(b-y)^\top A\,(b-y).
\]
The vector $\delta A_t^{-1}c_t$ is the Newton step: the accumulated derivative divided by
the accumulated second derivative. This is is so the directions where the logarithm has
curved a lot get small steps. The result may not be a valid portfolio (it could have
negative numbers or not add up to one), so the last operation $\Pi$ projects it, which
essentially means that it finds the valid portfolio that is closest to the solution. 
The distance is measured with $A_t$ so that differences in the directions where the logarithm 
curves the most count more than in the others. The values used are suggested by \textcite{agarwal06}
and used in the experiments of the book, $\beta=1$ and $\delta=1/8$. They are not tuned, because 
of the computational cost of ONS. The regret of ONS is $O(m^{1.5}\log(mn))$ \parencite{agarwal06} 
which grows very slowly with $n$, which means that ONS is universal.

\section{Comparison}

The three algorithms follow the winners, but they do it in different ways. UP uses the
whole history through the average of all the CRPs, EG only uses the last day, and ONS uses
the whole history with a Newton step.

In regret, UP is the best one: its regret is $O(m\log n)$, against $O(m^{1.5}\log(mn))$
for ONS and $O(\sqrt{n\log m})$ for EG. The bound of EG is the weakest, because it grows
with the square root of $n$, and it only holds if $\eta$ decreases with the number of days,
which is not the case in our experiments. In computation, EG is the cheapest, because its
update only needs the weights of the day before and the prices of the last day. ONS has to
solve a matrix problem every day, and UP is too expensive when there are many assets, which
is why we used the sampled version, and this version loses the guarantee.

All of them are compared with the BCRP, which is a fixed portfolio for the whole period. If
the best weights change over time, being close to the BCRP does not say much about how much
an algorithm earns, so it makes sense to look at other families of algorithms.

\section{Implementation}

The three classes follow the interface of Chapter~2. None of them remembers anything from
one day to the next: each call to \texttt{update} receives the whole history and computes
$b_{t+1}$ from it.

\texttt{EG} starts from the uniform portfolio and applies the update once for each day of
the history. \texttt{UP} draws the $K$ portfolios once, with a fixed seed, so it always
gives the same result. \texttt{ONS} goes through the history updating $A_t$ and $c_t$, and
solves the projection as a quadratic program with \texttt{cvxpy}
\parencite{diamond16,agrawal18}. The solver sometimes
returns tiny negative weights, around $10^{-9}$, where the result should be exactly zero,
so we set them to zero before the portfolio is validated.

Computing everything from the first day at each call is cheap for EG and UP, but not for
ONS. The projection has to be solved once for every day of the history, so for a run of
$n$ days the number of quadratic programs grows as $n^2/2$. This is why ONS is by far the
slowest strategy of the project.
